\documentclass[addpoints]{exam}
% leqno 將方程式編號放在左側

%\printanswers

\usepackage[top=2.5cm,bottom=2cm,left=2.5cm,right=2.5cm,headsep=10pt,a4paper]{geometry} % 引用頁面幾何套件，設定頁面邊界

\usepackage{siunitx} % 角度
% \input{NewCommands}

\usepackage{amsmath,amssymb,amsthm} % 引用常見的數學符號與設定


\theoremstyle{definition}
\newtheorem*{definition}{Definition}
\newtheorem*{theorem}{Theorem}
\newtheorem*{remark}{Remark}
\newcommand{\R}{\mathbb{R}}

\usepackage{lastpage} % 取得最後頁碼




\usepackage{graphicx} % 引用圖檔內嵌入套件
\graphicspath{{Pictures/}} % 設定圖檔目錄
\usepackage{xcolor} %引用顏色套件
\usepackage{enumerate}
\usepackage{float}
\usepackage{hyperref}
\usepackage{mdwlist} % use suspend & resume to successive enumerate

\usepackage{CJKutf8} % 設定中文

\def \lflr{\left\lfloor} %自定義 floor function
\def \rflr{\right\rfloor} %自定義

% question separation
\renewcommand{\questionshook}{%
  \setlength{\itemsep}{0.5cm}%
}

\allowdisplaybreaks %equation allowed to next page

\firstpageheader{\today}
                {}
                {}
\footer{}{\sf Page \thepage~of~\pageref{LastPage}}{}

\begin{document}

\begin{CJK*}{UTF8}{gbsn}
% Title
\begin{center}
    {\Large{\bf Introduction to Mathematical Analysis II\\
    Homework 9  Due  May 16  (Saturday), 2026\\
    Please submit your group homework online in PDF format.\\
You are expected to work collaboratively within your group. 
Please do not obtain complete solutions directly from AI tools. 
The purpose of this assignment is to develop your own mathematical reasoning and problem-solving skills. %  Homework X Brief Solution
  }} \\ 
    \large
    %\text{Due date: 3/9/2023}
\end{center}

\noindent\rule{16.2cm}{0.4pt}


\begin{questions}


% ============ Question x ============
\question  (20 pts) 



Here is a trick question: ``Are there any functions for which the Riemann
integral converges but the Lebesgue integral diverges?'' 
Show, however, that the improper Riemann integral
\[
\int_0^1 f(x)\,dx
\]
of
\[
f(x)=
\begin{cases}
\dfrac{\pi}{x}\sin\dfrac{\pi}{x}, & \text{if } x\neq 0,\\[1.2ex]
0, & \text{if } x=0,
\end{cases}
\]
exists and is finite, while the Lebesgue integral is infinite.

\medskip

\noindent
\textbf{Hint.}
Integration by parts gives
\[
\int_a^1 \frac{\pi}{x}\sin\frac{\pi}{x}\,dx
=
\left. x\cos\frac{\pi}{x}\right|_a^1
-
\int_a^1 \cos\frac{\pi}{x}\,dx.
\]
Why does this converge to a limit as \(a\to 0^+\)?

To check divergence of the Lebesgue integral, consider intervals
\[
\left[\frac{1}{k+1},\frac{1}{k}\right].
\]
On such an interval the sine of \(\pi/x\) is everywhere positive or everywhere
negative. The cosine is \(+1\) at one endpoint and \(-1\) at the other. Now use
the integration by parts formula again and the fact that the harmonic series
diverges.


\begin{solution}
\end{solution}

\question (20 pts) 
Assume that \(f\) is measurable on \(\mathbb{R}^2\), and assume that at least
one of the two iterated integrals
\[
\int_{\mathbb{R}}
\left[
\int_{\mathbb{R}} |f(x,y)|\,dx
\right]dy
\qquad\text{or}\qquad
\int_{\mathbb{R}}
\left[
\int_{\mathbb{R}} |f(x,y)|\,dy
\right]dx
\]
exists. Prove the following results:
\begin{enumerate}
    \item[\textup{a)}] \(f\in L(\mathbb{R}^2)\).

    \item[\textup{b)}]
    \[
    \iint_{\mathbb{R}^2} f
    =
    \int_{\mathbb{R}}
    \left[
    \int_{\mathbb{R}} f(x,y)\,dx
    \right]dy
    =
    \int_{\mathbb{R}}
    \left[
    \int_{\mathbb{R}} f(x,y)\,dy
    \right]dx.
    \]
\end{enumerate}


\begin{solution}


\end{solution}



\question (20 pts)
Let
\[
f(x,y)=e^{-xy}\sin x\sin y
\]
if \(x\ge 0\), \(y\ge 0\), and let \(f(x,y)=0\) otherwise. Prove that both
iterated integrals
\[
\int_{\mathbb R}
\left[
\int_{\mathbb R} f(x,y)\,dx
\right]dy
\qquad\text{and}\qquad
\int_{\mathbb R}
\left[
\int_{\mathbb R} f(x,y)\,dy
\right]dx
\]
exist and are equal, but that the double integral of \(f\) over \(\mathbb R^2\)
does not exist. Also, explain why this does not contradict the result in previous exercise.
\begin{solution} 

\end{solution}

\question ( 20 pts) 
Suppose that \(\mathcal E\subset C^0([a,b], \R )$ (
 the set of continuous functions $[a,b] \to \R$) is equicontinuous and bounded.
\begin{enumerate}
    \item[\textup{(a)}] Prove that
    \[
    \sup\{f(x):f\in\mathcal E\}
    \]
    is a continuous function of \(x\).

    \item[\textup{(b)}] Show that \textup{(a)} fails without equicontinuity.

    \item[\textup{(c)}] Show that this continuous-sup property does not imply
    equicontinuity.

    \item[\textup{(d)}] Assume that the continuous-sup property is true for each
    subset \(\mathcal F\subset\mathcal E\). Is \(\mathcal E\) equicontinuous?
    Give a proof or counterexample.
\end{enumerate}


\begin{solution} 
 
 \end{solution}
 
 
\question (10)
Is the sequence of functions \(f_n:\mathbb R\to\mathbb R\) defined by
\[
f_n(x)
=
\cos(n+x)
+
\log\left(
1+\frac{1}{\sqrt{n+2}}\sin^2(n^n x)
\right)
\]
equicontinuous? Prove or disprove.





\begin{solution} 
 
 \end{solution}



\question (10 pts)
If \(f:\mathbb R\to\mathbb R\) is continuous and the sequence
\[
f_n(x)=f(nx)
\]
is equicontinuous, what can be said about \(f\)?


   
\begin{solution} 
 
 \end{solution}

\

\end{questions}
\end{CJK*}
\end{document}

